\documentclass[11pt,a4paper]{article}

\usepackage[T1]{fontenc}
\usepackage[utf8]{inputenc}
\usepackage[cyr]{aeguill}
\usepackage[francais]{babel}

% %\renewcommand{\baselinestretch}{1.5}  % line spacing : 1,5
% \setlength{\oddsidemargin}{2,5cm}	% Marge gauche sur pages impaires
% \setlength{\evensidemargin}{2,5cm}	% Marge gauche sur pages paires
% \setlength{\textwidth}{16cm}	% Largeur de la zone de texte (17cm)
% \setlength{\topmargin}{2,5cm}	% Pas de marge en haut

% Les packages
%\usepackage[french]{babel}
\usepackage{fancybox}
\usepackage{graphics}
\usepackage{color}
\usepackage{latexsym}
\usepackage{amsmath,amsthm,amssymb}
\usepackage[plain]{fullpage} 
\usepackage{verbatim}
\usepackage{times,epsfig}
\usepackage{listings}
\usepackage{multirow}
\usepackage{amsmath,amsthm,amssymb,amscd,txfonts,bbm} % RAJOUT ICI !!!
\usepackage{graphics,epsfig} % RAJOUT ICI !!!
%\usepackage{algorithmic}
\usepackage[final]{pdfpages} 
\usepackage{fancyvrb}
\usepackage{enumerate}
\usepackage{listingsutf8}
\usepackage{tikz}
\usetikzlibrary{arrows,shapes,snakes,automata,backgrounds,petri,positioning}
\usepackage{tikz-qtree}
\usepackage{comment}


\newenvironment{qv}
{\quote\Verbatim}
{\endquote\endVerbatim}

\usepackage{lastpage}

\newcommand{\terme}[1]{\texttt{#1}}
\newcommand{\termSet}[1]{\{\texttt{#1}\}}
\newcommand{\guill}[1]{<<\,#1\,>>}
% Macro pour l'inclusion de figure  
\newcommand{\fig}[2]{%
  \begin{figure*}[hbt]
   \begin{center}
      \includegraphics[width=6cm]{#1}
  \caption{\label{#1}  #2}
  \end{center}
 \end{figure*}
}
\newcommand{\figTex}[2]{%
  \begin{figure*}[hbt]
 \centerline{\input{#1.pstex_t}}
  \caption{\label{#1}  #2}
 \end{figure*}
}

% Macro commentaires
\definecolor{turquoise}{rgb}{0.20 ,0.60, 0.60}
\definecolor{vert}{rgb}{0.40,  0.60 ,0.00}
\definecolor{violet}{rgb}{0.80 , 0.00 , 0.80}
\definecolor{bleu}{rgb}{0.20,  0.20,  0.80}
\definecolor{rose}{rgb}{1.00,  0.20,  0.40}
\def\remVG#1{\textcolor{vert}{\texttt{Virginie: #1}}}
\def\comVG#1{\begin{footnotesize}\textcolor{vert}{\texttt{Precision de Virginie: #1}}\end{footnotesize}}

\lstset{basicstyle=\footnotesize,showstringspaces=false,language=XML,breaklines=false,columns=fullflexible,frame=shadowbox, captionpos=b}

\newenvironment{solution}[0]{
~\\ \color{red} \textbf{Solution :}\color{black}}{%
}
%\excludecomment{solution}


\parindent=0pt           
                             
\usepackage{url}

\newtheorem{Exercice}{Definition}

\lstset{language=XML,basicstyle=\footnotesize,
numberstyle=\footnotesize, 
breaklines=true
showspaces=false,         
showstringspaces=false,
showtabs=false,
frame=single,
tabsize=4}   


\begin{document}

\title{Examen NFE204 - Session 1}
\date{03 février 2015}

\maketitle

\begin{center}
{\bf \Large Consignes}\\
\begin{tabular}{| p{15cm} |}
\hline
\\
Dur\'ee de l'examen : 3h00\\
Documents non autoris\'es;  Calculatrice autoris\'ee\\

Les exercices sont ind\'ependants les uns des autres. Merci de soigner votre \'ecriture.\\
Ce sujet contient 2 pages.\\


\textit{Important}\,: nous n'évaluons pas les réponses par la syntaxe. Ne vous
inquiétez pas d'utiliser un point-virgule à la place d'un point d'exclamation ou
autres détails mineurs
\\
\hline
\end{tabular}
\end{center}

\section*{Première partie : recherche d'information (8 pts)}

Voici quelques extraits d'un discours politique célèbre (un peu modifié pour les besoins de la cause). Chaque
extrait correspond à un document, numéroté a$_i$.

\begin{small}
\begin{verbatim}
(a1) Moi, président de la République, je ne serai pas le chef de la majorité, 
je ne recevrai pas les parlementaires de la majorité à l’Elysée.
\end{verbatim}
\end{small}

\begin{small}
\begin{verbatim}
(a2) Moi, président de la République, je ne traiterai pas mon Premier ministre de collaborateur.
\end{verbatim}
\end{small}

\begin{small}
\begin{verbatim}
(a3) Moi, président de la République,  les ministres de la majorité
ne pourraient pas cumuler  leurs fonctions avec un mandat parlementaire ou local.
\end{verbatim}
\end{small}


\begin{small}
\begin{verbatim}
(a4) Moi, président de la République, il y aura un code de déontologie  pour les ministres 
et parlementaires qui ne pourraient pas rentrer dans un conflit d’intérêt.
\end{verbatim}
\end{small}

\begin{enumerate}
  \item Rappeler la notion de \textit{stop word} (ou "mot vide") et donner la liste
         de ceux que vous choisiriez dans les textes ci-dessus.
         
         
\begin{solution}
Mot très fréquents et peu porteurs de sens pour l'interprétation d'un contenu: le, de, la,
ne, pas, un, etc.
\end{solution}

  \item Outre ces mots vides, pouvez-vous identifier certains mots dont l'idf tend vers 0 (en appliquant le logarithme)? Lesquels?  
  
           
\begin{solution}
Moi, président, république: cas particuliers pour ce corpus, il apparaissent dans chaque document.
\end{solution}

  \item Présentez la matrice d'incidence pour le vocabulaire suivant:  ''majorité, ministre, déontologie, parlementaire''.
  Vous indiquerez  l'idf pour chaque terme (sans logarithme), et le tf pour chaque paire (terme, document). 


\begin{solution}
  
\begin{center}
\begin{tabular}{l|c|c|c|c}
       & \textbf{a1} & \textbf{a2} & \textbf{a3} &\textbf{a4} \\ \hline  
   majorité (2) & 2 & 0 & 1 & 0 \\
   ministre (4/3) & 0 & 1 &  1 & 1\\
   déontologie (4) & 0 & 0 & 0 & 1\\
   parlementaire (4/3) & 1 & 0 & 1 & 1 \\ \hline
   \end{tabular}
\end{center}

\end{solution}

  \item Donner les résultats classés par similarité cosinus basée sur les tf (on ignore l'idf) pour les requêtes suivantes.
  
  \begin{itemize}
    \item majorité; expliquez le classement;
    \item ministre; expliquez le classement du premier document;
    \item déontologie et ministre;  qu'est-ce
         qui changerait si on prenait en compte l'idf?
    \item majorité et ministre; qu'obtiendrait-t-on avec une requête Booléenne? Commentaire?
  \end{itemize}
\begin{solution}
Les normes
\begin{itemize}
  \item $||a1|| = \sqrt{4+ 1}$; $||a2|| = \sqrt{1}$; $||a3|| = \sqrt{3}$; $||a4|| = \sqrt{3}$ 
\end{itemize}

Les cosinus (requête non normalisée).

\begin{itemize}
  \item majorité: a1 : $\frac{2}{\sqrt{5}}$;  a3 : $\frac{1}{\sqrt{3}}$; a1 est premier car majorité apparaît 2 fois.
  
  \item ministre: a2 : $\frac{1}{\sqrt{1}}$;  a3 : $\frac{1}{\sqrt{3}}$,  a4 : $\frac{1}{\sqrt{3}}$; le document a2 ne
      parle que de ministre, et rien d'autre.
  
  \item déontologie et ministre: a2 : $\frac{1}{\sqrt{1}}$;  a3: $\frac{1}{\sqrt{3}}$, 
          a4 : $\frac{2}{\sqrt{3}}$;
     déontologie ayant un idf supérieur à ministre, a3 serait classé avant a2, 
  
  \item majorité et ministre: a1 : $\frac{2}{\sqrt{5}}$;  a2 : $\frac{1}{\sqrt{1}}$; a3 : $\frac{2}{\sqrt{3}}$; 
  a4 : $\frac{1}{\sqrt{3}}$; avec une requête Booléenne, a1 n'apparaîtrait pas du tout!
    
\end{itemize}

\end{solution}  

   \item Calculez la similarité entre a3 et a4; puis entre a3 et a1. Qui est le plus proche de a3?
   
\begin{solution}
Il s'agit toujours de la similarité cosinus. La similarité par fonction Euclidienne
sur les vecteurs \emph{normalisés} fonctionne aussi.

\begin{itemize}
  \item cos(a3, a4) = $\frac{1 +  1}{\sqrt{3} \times \sqrt{3}}$;
  \item cos(a3, a1) = $\frac{2 +  1}{\sqrt{5} \times \sqrt{3}}$.
\end{itemize}
a1 est plus proche de a3. 

\end{solution}
  \end{enumerate}

\section*{Seconde partie : Pig et MapReduce (6 pts)}

Un système d'observation spatial capte des signaux en provenance de planètes
situées dans de lointaines galaxies.  Ces signaux sont stockés dans une collection \textit{Signaux} de la forme

\begin{center}
  Signaux (idPlanète, date, contenu)
\end{center}

Le but est de déterminer si ces signaux peuvent être émis par une intelligence extra-terrestre.
Pour cela les scientifiques ont mis au point les fonctions suivantes:

  \begin{itemize}
     \item \emph{Fonction de structure}: $f_S(c): Bool$, prend un contenu en entrée, et renvoie \texttt{true}
        si le contenu présente une certaine structure, \texttt{false} sinon.
     \item \emph{Fonction de détecteur d'Aliens}: $f_D(<c>): real$, prend une liste de contenus \emph{structurés} en entrée, et renvoie 
        un indicateur entre 0 et 1 indiquant la probabilité que ces contenus soient écrits en  langage extra-terrrestre,
        et donc la \emph{présence d'Aliens}!
  \end{itemize}
  
Bien entendu, il y a beaucoup de signaux: c'est du Big Data.

\medskip
\textbf{Questions}.

\begin{enumerate}
  \item Ecrire un programme Pig latin qui produit, pour chaque planète, l'indicateur de présence d'Aliens
     par analyse des contenus provenant de la planète.
  \item Donnez un programme MapReduce qui permettrait d'exécuter ce programme Pig en distribué
      (indiquez la fonction de Map, la fonction de Reduce, dans le langage ou pseudo-code qui vous convient).
  \item Ecrire un programme Pig latin qui produit, pour chaque planète et pour chaque jour, 
     le rapport entre  contenus structurés et non structurés reçus de cette planète. 
\end{enumerate}


\begin{solution}
\begin{enumerate}
 \item À la syntaxe près:
\begin{verbatim}
// On ne prend que les signaux avec contenu structuré
structurés = filter Signaux by fS(contenu) = true;
// On le regroupe par planète
signauxPlanete = group structurés  by idPlanete;
// Maintenant, dans signauxPlanete, on a un liste de contenus structurés
// par planète: on applique le détecteur à cette liste
présenceAliens = foreach signauxPlanète generate group, 
                   fD(signauxPlanete) as indicateur;
\end{verbatim}

\item La fonction de  Map émet les signaux pour lesquels $f_S(s)=true$. La clé
   des paires intermédiaires est celle sur laquelle on veut effectuer
   le regroupement, soit \texttt{idPlanète}.

\begin{verbatim}
function map(signal)
{
  if (fS(signal.contenu) == true) {
     emit (signal.idPlanète, signal.contenu)
  }
}
\end{verbatim}

La fonction de Reduce reçoit l'id d'une planète, et la liste
des contenus structurés. On applique fD à cette liste.
\begin{verbatim}
function reduce(idPlanèe, [contenus])
{
  return (idPlanète, fD(contenus);
}
\end{verbatim}

\item Un poil plus compliqué: on constitue la liste des contenus structurés,
   la liste des contenus non structurés, on les regroupe (\texttt{cogroup} va
   très bien pour ça) par planète et date, et on calcule le ratio.
   
\begin{verbatim}
structurés = filter Signaux by fS(contenu) == true;
nonstructurés = filter Signaux by fS(contenu) == false;
signauxClassés = cogroup structurés by (idPlanete,date), 
                         nonstructurés by (idPlanete, date);
ratioStructurés = foreach signauxClassés 
                  generate group, count(structurés) / count(nonstructurés) ;
\end{verbatim}
\end{enumerate}
\end{solution}


\section*{Troisième partie : questions de cours (6 pts)}


\begin{enumerate}

  \item En recherche d'information, qu'est-ce que le rappel? qu'est-ce que la précision?
  
  \item Deux techniques fondamentales vues en cours sont la réplication et le partitionnement.
     Rappelez brièvement leur définition, et indiquez leurs rôles respectifs. Sont-elles complémentaires? Redondantes?
       
  \item Qu'est-ce qu'une architecture multi-nœuds, quels sont ses avantages et inconvénients?
        
  \item Vous avez 500 TOs de données, et vous pouvez acheter des serveurs pour votre \textit{cloud}
        avec chacun 32 GO de mémoire et 
      10 TOs de disque. Le coût unitaire d'un serveur eest de 500 Euros.
      
      Quelle est la configuration de votre grappe de serveurs la moins coûteuse (financièrement) et combien de temps prend
      au minimum la lecture complète de la collection avec cette solution?
      
      \begin{solution}
      Il faut au minimum 50 serveurs, soit 25 000 Euros. La lecture en en parallèle 
      de disques de 10 TO prend (à 100 MO/s) plus de 20 heures.
      
      On peut ajouter la réplication (facter 2 ou 3) au calcul.
      \end{solution}
      
   \item Même situation:   quelle est la configuration qui assure le maximum d'efficacité, et quel est son coût financier?
       
      \begin{solution}
      L'idéal est de tout stocker en RAM. mais là il faut $\frac{500 000}{32} = 15 625$ serveurs
      au minimum.
      \end{solution}
       
    \item Rappelez le principe de l'éclatement d'un fragment dans le partitionnement par intervalle.
\end{enumerate}

\end{document}
